\section{A Continuing Case: Bayes Institute's Starting Position}

\subsection{The Reality of Day Zero}

The general theory becomes clearer when applied to a real founder. The following case is based on the founding process for Bayes Institute. It is \textbf{not} an independent audit; claims about background are founder-reported, and commercial demand remains unverified.

The long-term ambition is substantial: launch a lean STEM education venture and, over decades, evolve into a research-active university with world-class faculty, physical campuses, and international scientific standing. The immediate starting reality is different: a single founder with a full-time technology job, extensive software and teaching experience, a laptop, internet connection, limited operating capital, and zero alumni brand or campus.

\subsection{What the Founder Plausibly Brings}

\begin{tcolorbox}[
  enhanced,
  colback=white,
  colframe=BayesBorder!80!white,
  boxrule=0.4pt,
  arc=2pt,
  left=3pt,right=3pt,top=3pt,bottom=3pt,
  before skip=3pt,after skip=5pt
]
\scriptsize
\setlength{\tabcolsep}{2pt}
\renewcommand{\arraystretch}{1.18}
\begin{tabularx}{\linewidth}{@{} >{\bfseries\color{BayesInk}}p{1.75cm} >{\color{BayesInk}}p{1.95cm} >{\color{BayesTeal}}p{1.95cm} >{\color{BayesSecondary}}X @{}}
\toprule
\textbf{Dimension} & \textbf{Starting Evidence} & \textbf{What It Helps With} & \textbf{What It Does Not Prove} \\
\midrule
Technical depth & AI/ML, distributed systems, statistical math & Designing rigorous labs and evaluation sets & Ability to run a large multi-faculty institute. \\
Pedagogical skill & Extensive prior curriculum design and teaching & Diagnosing misconceptions; building rubrics & Commercial willingness to pay for a new brand. \\
Industry view & Enterprise software and edtech experience & Identifying gaps in current university curricula & That personal observations represent whole market. \\
Quantitative rigor & Advanced formal study in statistics & Teaching foundations and probabilistic betting & Demand for a math-heavy premium program. \\
Execution & Small training deliveries and direct sales & Running lean pilots and course cohorts & Repeatable institutional enterprise contracts. \\
Tooling, systems & Laptop, workstation, open software & Rapid low-cost prototyping and remote cohorts & Operating capability for a physical campus. \\
\bottomrule
\end{tabularx}
\end{tcolorbox}

\textbf{The most plausible starting advantage for Bayes Institute is the intersection of technically deep instruction, AI systems capability, statistical foundations, and familiarity with technical curriculum operations.}

\subsection{What Is Missing or Constrained}

There is no basis yet to assume an automatic pipeline of paying students, measured employment outcomes, employer recognition of Bayes certificates, or university partnership contracts. These are future assets that must be earned.

Furthermore, the founder intends to maintain full-time employment during early validation. This provides essential household financial stability but severely limits daytime responsiveness. All IP, customer communication, and outside engagements must strictly honor employment policies. Finally, unverified contact lists from past jobs cannot be commercialized without explicit permission.

\subsection{3 Possible Starting Routes}

\textbf{Route 1: Narrow D2C technical cohort.} A small program for working software engineers needing deep AI evaluation and statistical systems skills. Leverages founder teaching; tests learner willingness to pay.

\textbf{Route 2: College applied-course pilot.} A bounded curriculum pilot integrated into an engineering college. Leverages pedagogical credentials; tests institutional procurement and renewal.

\textbf{Route 3: Enterprise AI capability pilot.} A specialized corporate workshop paired with a workflow diagnostic. High potential revenue; high conflict risk with daytime hours.

\begin{tcolorbox}[
  enhanced,
  colback=white,
  colframe=BayesBorder!80!white,
  boxrule=0.4pt,
  arc=2pt,
  left=3pt,right=3pt,top=3pt,bottom=3pt,
  before skip=3pt,after skip=5pt
]
\scriptsize
\setlength{\tabcolsep}{2pt}
\renewcommand{\arraystretch}{1.18}
\begin{tabularx}{\linewidth}{@{} >{\bfseries\color{BayesInk}}p{1.8cm} >{\color{BayesInk}}p{1.9cm} >{\color{BayesInk}}p{1.9cm} >{\color{BayesSecondary}}X @{}}
\toprule
\textbf{Fit Dimension} & \textbf{D2C Cohort} & \textbf{College Pilot} & \textbf{Enterprise Pilot} \\
\midrule
Delivery fit & Strong (evenings) & Moderate & Weak (daytime) \\
Buyer access & Untested & Emerging & Unproven \\
Brand trust & Moderate & Institutional & High risk \\
Chief test & Will devs pay? & Will colleges renew? & Will CIOs sign? \\
\bottomrule
\end{tabularx}
\end{tcolorbox}

\subsection{Claims Bayes Must Resist Making}

It would be tempting to market the founder's background as proof that Bayes will immediately transform careers or guarantee hiring. That would exceed the evidence. At day zero, the legitimate claim is narrower: the founder has verified competence and can demonstrate the curriculum and evaluation rigor before buyers commit.

\textbf{Distinguish a founder's opening advantage from an institution's mature moat: founder competence can open doors, but systems, faculty, alumni outcomes, and research reputation must be compounded over years.}

The next operational tests for Bayes are clear:
\begin{enumerate}[label=\textbf{\arabic*.},leftmargin=11pt,itemsep=1.5pt,topsep=1.5pt,parsep=0pt]
  \item Build an unseen technical diagnostic and invite external experts to critique it.
  \item Run 10 discovery interviews with working developers and 5 with college deans.
  \item Offer a bounded paid pilot with milestone refund clauses.
  \item Simulate 2 weeks of teaching and grading under full-time job conditions.
\end{enumerate}
